% Recuperación literal de nucleo.tex REV08, líneas 10--32.
% Ampliación específica de II, situada después del núcleo común inalterado.
\subsection{Correspondance entre support positionnel et calendrier orienté}
\label{ii:ampliacion-posicional-calendario}

La représentation du calendrier distingue les positions d'APP et les
transitions orientées de son évolution. Nous reprenons cette correspondance pour
expliciter la lecture du segment contenant le raccordement entre tours.

Considérons l'alphabet à neuf indices
\[
 D_9=\{1,2,3,4,5,6,7,8,9\}.
\]
Sa disposition horizontale et verticale produit les quatre-vingt-une positions
de \(D_9^2\). Pour décrire le mouvement, on utilise également le cycle
orienté \(C_9\). Ses neuf arêtes reçoivent, dans l'ordre, les indices de
\(D_9\). L'indice d'une arête spécifie une transition entre états,
tandis qu'un sommet spécifie sa frontière. L'application qui associe
à chaque arête son indice permet d'utiliser le même alphabet dans le support
positionnel et dans le calendrier ; elle conserve la distinction entre position,
intervalle orienté et état transporté.

Dans le relèvement entier du calendrier, le premier tour comprend
les neuf intervalles consécutifs \([0,1],\ldots,[8,9]\). Le prolongement
\([9,10]\) commence le tour suivant avec l'indice un et une mémoire
mise à jour. Le segment représenté jusqu'à dix montre les deux tours à leur
point de raccordement. Cette représentation spécifie les neuf transitions d'un
tour et la transition suivante ; l'holonomie de l'état se construit
ensuite par la composition effective du TPK.

\begin{figure}[htbp]
\centering
\begin{tikzpicture}[x=1.05cm,y=1cm,>=Latex,font=\small]
 \draw[very thick,ink] (0,0)--(9,0);
 \draw[very thick,dashed,accent] (9,0)--(10,0);
 \foreach \x in {0,...,10}{
   \draw (\x,-.08)--(\x,.08);
   \node[below=5pt,font=\scriptsize] at (\x,0) {$\x$};
 }
 \foreach \x in {1,...,9}{
  \draw[->,ink] (\x-.85,.35)--(\x-.15,.35);
  \node[above=3pt] at (\x-.5,.35) {$\x$};
 }
 \draw[->,accent] (9.15,.35)--(9.85,.35);
 \node[above=3pt,text=accent] at (9.5,.35) {$1$};
 \draw[decorate,decoration={brace,amplitude=5pt},ink] (0,1)--(9,1)
  node[midway,above=7pt] {Neuf intervalles orientés : un tour};
 \node[align=center] at (5,-1.05)
  {Indice de transition : $g(t)=1+(t\bmod9)$\\
   Mémoire du calendrier relevé : $m(t)=\lfloor t/9\rfloor$};
 \node[align=center] at (5,-1.85)
  {$g(t+9)=g(t),\qquad m(t+9)=m(t)+1$};
\end{tikzpicture}
\caption{Relèvement du cycle de neuf transitions. Les nombres supérieurs
étiquettent les intervalles orientés ; les nombres inférieurs situent leurs frontières.
L'intervalle discontinu montre la première transition du tour suivant.
L'indice de phase revient et le quotient du calendrier augmente. Ce schéma
représente le calendrier ; l'action sur les résidus, quotients, orientation,
frontière et mémoire est définie par les opérations du TPK.}
\label{fig:ii-segmento-generador}
\end{figure}

